\section{The Content Account Does Not Survive Re-analysis}
\label{sec:reanalysis}

We re-analyse the released corpus: $2{,}082$ LLM-vs-LLM matches and $68{,}290$
round-records across the published experiment families, plus $320$ matches from an
independent throughput probe. Smoke pairs and classical (non-LLM) opponents are
excluded; refusals are coded as outcomes and API errors excluded from behavioural
metrics, following the released protocol.

\subsection{Punishment rules are rare and predict nothing}
\label{sec:rules}

We detect contingent punishment rules with a high-precision detector requiring an
\emph{adverse} consequence co-occurring with a \emph{peer-deviation} trigger in the
same sentence, with cooperative-reciprocity phrasings vetoed (\S\ref{sec:methods}
documents its construction and a retraction). Across $540$ free-text IPD matches:

\begin{itemize}[leftmargin=1.1em,itemsep=2pt,topsep=2pt]
\item rules appear in $17.2\%$ of matches and $0.6\%$ of messages;
\item they are \emph{not} elevated after a defection: round-level hazard difference
$+0.002$ $[-0.004,+0.008]$;
\item they do not aid recovery. Restricted to rounds in which a defection just
occurred, $P(\text{mutual cooperation next round})$ is $0.164$ with a rule versus
$0.116$ without: $+0.048$ $[-0.036,+0.142]$;
\item they do not predict the match outcome: IPD lock-in $-0.099$ $[-0.211,+0.012]$
($n{=}93$ vs $447$); Bertrand $K$ $+0.143$ $[-0.103,+0.517]$ ($n{=}16$ vs $544$, too
thin to interpret).
\end{itemize}

\noindent The grim-trigger scheme quoted in the prior paper is genuinely present in
the transcripts. It is not doing the work.

\subsection{No content class predicts the outcome}
\label{sec:horserace}

Observational mediation is vulnerable to reverse causation---agents may threaten
\emph{because} a match is going badly. We therefore score every candidate mediator on
the \emph{opening} rounds and predict cooperation in the \emph{closing} half,
controlling for how the opening actually went. A mediator measured before the outcome
window cannot be a response to it. Table~\ref{tab:horserace} gives the result.

\begin{table}[t]
\centering
\caption{\textbf{Mediator horse-race, temporal design.} Content classes scored on
rounds $0$--$2$; outcome is closing-half cooperation. IPD free-text, $n{=}484$.
Everything the agents \emph{say} is null; what predicts the ending is what they
\emph{did} at the start.}
\label{tab:horserace}
\small
\begin{tabularx}{0.8\textwidth}{X r l}
\toprule
term & coefficient & 95\% CI \\
\midrule
opening threat            & $+0.055$ & $[-0.005, +0.117]$ \\
opening joint-plan        & $-0.048$ & $[-0.119, +0.023]$ \\
opening affirmation       & $+0.102$ & $[-0.023, +0.237]$ \\
opening solicitation      & $-0.016$ & $[-0.071, +0.041]$ \\
opening message length (z)& $-0.017$ & $[-0.054, +0.016]$ \\
\textbf{opening cooperation (z)} & $\mathbf{+0.228}$ & $\mathbf{[+0.196, +0.255]}$ \\
\bottomrule
\end{tabularx}
\end{table}

Message length is null here and trends negative within sustained matches
(opening-length tertiles: short $0.827$, mid $0.795$, long $0.775$). This is the
opposite of what a bandwidth account predicts.

\subsection{The channel both displaces the opening and slows the decay}
\label{sec:decomp}

Decomposing by rung (Table~\ref{tab:decomp}), the channel raises opening cooperation
from $0.622$ to $0.815$ \emph{and} slows the subsequent decay. Because decay is
floor-confounded---an arm sitting near zero cannot fall---we re-test within matches
that opened above $C>0.8$, where every arm has equal room to fall. The channel effect
on persistence survives: $\mathrm{L2}-\mathrm{L0}$ decay $+0.352$ $[+0.292,+0.411]$.
Both components are real, which is what makes the defense in \S\ref{sec:p4} an
apportionment question rather than a foregone conclusion.

\begin{figure}[t]
\centering
\includegraphics[width=0.85\textwidth]{figures/f2_trajectory.pdf}
\caption{\textbf{Cooperation by round and rung (IPD).} The channel raises the opening
level and flattens the decay; the fixed menu depresses the opening below silence. The
shaded band is the opening window targeted by the defense in \S\ref{sec:p4}. Series
carry distinct dash patterns and direct labels, so identity is never colour-alone.}
\label{fig:traj}
\end{figure}

\begin{table}[t]
\centering
\caption{\textbf{Opening displacement and persistence, IPD.} Opening $=$ rounds
$0$--$2$; closing $=$ second half. Right-hand block restricts to matches opening at
$C>0.8$ so decay is not floor-confounded.}
\label{tab:decomp}
\small
\begin{tabular}{lrrrr|rrr}
\toprule
& \multicolumn{4}{c|}{all matches} & \multicolumn{3}{c}{opened at $C>0.8$} \\
rung & $n$ & open & close & decay & $n$ & close & decay \\
\midrule
L0 none      & 314 & 0.622 & 0.248 & $-0.373$ & 124 & 0.400 & $-0.549$ \\
L1 menu      & 115 & 0.301 & 0.069 & $-0.232$ &  14 & 0.374 & $-0.590$ \\
L2 free text & 366 & 0.815 & 0.644 & $-0.171$ & 267 & 0.791 & $-0.197$ \\
L3 private   & 118 & 0.820 & 0.734 & $-0.086$ &  85 & 0.825 & $-0.158$ \\
\bottomrule
\end{tabular}
\end{table}

\subsection{What does separate the arms}
\label{sec:predicate-obs}

Table~\ref{tab:predicate} scores every arm for whether it \emph{admits} a joint
proposal. The separation is complete, and it survives the obvious alternatives. The
\texttt{C1 restricted} arm is self-authored free text with a message every round,
which rules out authorship and channel presence; it carries own intentions and
behaves like the fixed menu.

\begin{figure}[t]
\centering
\includegraphics[width=0.86\textwidth]{figures/f1_predicate.pdf}
\caption{\textbf{The predicate separates every arm.} Cooperation lock-in by arm,
grouped by whether the channel admits a joint proposal. \texttt{C1 free text, own
intention only} is self-authored free text with a message every round and still sits at
$0.00$, which rules out authorship and channel presence as the active ingredient. At
the arm level the grouping is confounded with message length; \S\ref{sec:p1} breaks
that confound.}
\label{fig:predicate}
\end{figure}

\begin{table}[t]
\centering
\caption{\textbf{The predicate across seven arms.} Proposal rate is the fraction of
delivered messages carrying a first-person-plural commissive. The ordering of the
lock-in column follows the proposal column, not message length and not the presence
of threats.}
\label{tab:predicate}
\small
\begin{tabularx}{\textwidth}{X r r r r l}
\toprule
arm & prop.\ rate & int.-only & opening $C$ & lock-in & Wilson 95\% CI \\
\midrule
L0 none                       & 0.000 & 0.000 & 0.630 & 0.113 & [0.08, 0.15] \\
L1 menu                       & 0.000 & 1.000 & 0.318 & 0.057 & [0.03, 0.12] \\
C1 selected                   & 0.000 & 1.000 & 0.217 & 0.050 & [0.01, 0.24] \\
C1 restricted                 & 0.000 & 1.000 & 0.175 & 0.000 & [0.00, 0.16] \\
C1 free                       & 1.000 & 0.000 & 0.850 & 0.550 & [0.34, 0.74] \\
L2 observed                   & 0.918 & 0.022 & 0.845 & 0.599 & [0.55, 0.65] \\
L3 private                    & 0.939 & 0.013 & 0.825 & 0.651 & [0.56, 0.73] \\
\bottomrule
\end{tabularx}
\end{table}

At the arm level this is confounded with message length (\texttt{C1 restricted}
averages $20.7$ characters, \texttt{C1 free} $143$). The message-level evidence breaks
the confound in both directions---$172$ proposals under $40$ characters, $4{,}293$
own-intention messages of $20+$ characters, some over $200$---but a clean test needs
the manipulation, which is \S\ref{sec:p1}.

\subsection{Convergent evidence from an independent apparatus}
\label{sec:sms}

The throughput probe is a separate codebase in which each agent may send one
$\le\!60$-character SMS per round against a finite credit budget. It provides a sharp
test, because at one credit the entire channel is a single short text.

\begin{table}[t]
\centering
\caption{\textbf{The one-credit cell.} The single $60$-character text that lifts a
pair from $0.45$ to $0.80$ lock-in carries \emph{no} punishment rule. The pair that
fails to coordinate produces threats abundantly and still does not.}
\label{tab:sms}
\small
\begin{tabular}{lrrr}
\toprule
pair & credits & lock-in & rule rate \\
\midrule
qwen $\times$ deepseek & 0 & 0.45 & 0.00 \\
qwen $\times$ deepseek & \textbf{1} & \textbf{0.80} & \textbf{0.00} \\
qwen $\times$ llama    & 0 & 0.00 & 0.00 \\
qwen $\times$ llama    & 8 & 0.25 & 0.45 \\
\bottomrule
\end{tabular}
\end{table}

Every distinct text at one credit is a bare joint proposal:

\begin{excerpt}{Throughput probe, 1 credit, qwen $\times$ deepseek}
``Let's both HOLD for max profit.'' \quad ``Let's both HOLD for mutual benefit.''
\quad ``Let's HOLD, split the profit. $3{+}3 > 5{+}0$.''
\end{excerpt}

\noindent Threats are what these agents produce when coordination is \emph{not}
working---which is also why a naive mediation analysis assigns them a spurious
negative coefficient (\S\ref{sec:methods}).

% ======================================================================
